\documentclass{article}
\usepackage{amsmath,amssymb,amsthm}
\usepackage{algorithm}
\usepackage{algpseudocode}
\usepackage{enumitem}
\newtheorem{theorem}{Theorem}
\newtheorem{lemma}{Lemma}
\newtheorem{assumption}{Assumption}
\newcommand{\prox}{\operatorname{prox}}
\newcommand{\grad}{\nabla}
\newcommand{\R}{\mathbb{R}}
\newcommand{\E}{\mathbb{E}}

\begin{document}

\section*{Proximal Gradient Method (PGM) for Non‑Convex Composite Optimization}

\section*{Mathematical Formulation}
We consider the composite problem  

\[
\min_{x\in\R^{d}} \; \Phi(x)\quad\text{with}\quad 
\Phi(x):=f(x)+g(x),
\]

where  

\begin{itemize}[nosep]
\item $f:\R^{d}\to\R$ is \emph{smooth} (continuously differentiable) but possibly non‑convex,
\item $g:\R^{d}\to\R\cup\{+\infty\}$ is proper, lower semicontinuous and possibly non‑smooth (e.g. $\ell_{1}$‑regularizer).
\end{itemize}

For a stepsize $\eta>0$, the proximal operator of $g$ is  

\[
\prox_{\eta g}(v):=\arg\min_{u\in\R^{d}}\Big\{ g(u)+\frac{1}{2\eta}\|u-v\|^{2}\Big\}.
\]

The algorithm generates a sequence $\{x_{k}\}_{k\ge0}$ by  

\[
\boxed{\;x_{k+1}= \prox_{\eta g}\bigl(x_{k}-\eta\,\grad f(x_{k})\bigr)\; } \tag{1}
\]

with a fixed stepsize $\eta\in(0,1/L]$, where $L$ is the Lipschitz constant of $\grad f$ (see Assumption~\ref{ass:L}).

\section*{Pseudocode}
\begin{algorithm}[H]
\caption{Proximal Gradient Method (PGM)}
\begin{algorithmic}[1]
\Require Initial point $x_{0}\in\R^{d}$, stepsize $\eta\in(0,1/L]$,
        max. iterations $K$.
\For{$k=0,1,\dots,K-1$}
    \State $v_{k}\gets x_{k}-\eta\,\grad f(x_{k})$ \Comment{gradient step}
    \State $x_{k+1}\gets \prox_{\eta g}(v_{k})$ \Comment{proximal step}
\EndFor
\State \Return $x_{K}$
\end{algorithmic}
\end{algorithm}

\section*{Dry Run Example}
Take the scalar problem  

\[
f(w)=(w-3)^{2},\qquad g\equiv0,
\]
so $\Phi(w)=f(w)$.  
We use $w_{0}=0$, stepsize $\eta=0.1$, and run three iterations.

\begin{center}
\begin{tabular}{c|c|c|c}
$k$ & $\grad f(w_{k})$ & $w_{k+1}=w_{k}-\eta\grad f(w_{k})$ & $\Phi(w_{k+1})$\\ \hline
0 & $2(w_{0}-3)=-6$ & $0-0.1(-6)=0.6$ & $(0.6-3)^{2}=5.76$\\
1 & $2(0.6-3)=-4.8$ & $0.6-0.1(-4.8)=1.08$ & $(1.08-3)^{2}=3.6864$\\
2 & $2(1.08-3)=-3.84$ & $1.08-0.1(-3.84)=1.464$ & $(1.464-3)^{2}=2.3665$
\end{tabular}
\end{center}

The objective value decreases monotonically, illustrating the descent property of (1).

\newpage

\section*{Assumptions}
\begin{assumption}[Smoothness of $f$]\label{ass:L}
$f$ is $L$‑smooth, i.e. for all $x,y\in\R^{d}$  

\[
\|\grad f(x)-\grad f(y)\|\le L\|x-y\|.
\]
Equivalently,  

\[
f(y)\le f(x)+\langle\grad f(x),y-x\rangle+\frac{L}{2}\|y-x\|^{2}. \tag{2}
\]
\end{assumption}

\begin{assumption}[Bounded Below]\label{ass:bound}
The composite objective $\Phi$ is bounded below: $\Phi^{\inf}:=\inf_{x}\Phi(x)>-\infty$.
\end{assumption}

\begin{assumption}[Stepsize]\label{ass:step}
The stepsize satisfies $0<\eta\le 1/L$.
\end{assumption}

\section*{Main Convergence Theorem}
\begin{theorem}\label{thm:main}
Let $\{x_{k}\}$ be generated by \eqref{1} under Assumptions~\ref{ass:L}–\ref{ass:step}.  
Define the \emph{gradient mapping}  

\[
G_{\eta}(x):=\frac{1}{\eta}\bigl(x-\prox_{\eta g}(x-\eta\grad f(x))\bigr).
\]

Then for any integer $T\ge1$,

\[
\frac{1}{T}\sum_{k=0}^{T-1}\|G_{\eta}(x_{k})\|^{2}
\;\le\;
\frac{2}{\eta T}\bigl(\Phi(x_{0})-\Phi^{\inf}\bigr). \tag{3}
\]

Consequently, there exists at least one index $k\in\{0,\dots,T-1\}$ such that  

\[
\|G_{\eta}(x_{k})\|\le
\sqrt{\frac{2}{\eta T}\bigl(\Phi(x_{0})-\Phi^{\inf}\bigr)}=
O\!\left(\frac{1}{\sqrt{T}}\right). \tag{4}
\]

Thus the algorithm produces an \emph{approximate stationary point} in the sense  

\[
0\in\grad f(x_{k})+\partial g(x_{k})+ \mathcal{O}\!\bigl(\|G_{\eta}(x_{k})\|\bigr).
\]
\end{theorem}

\section*{Theoretical Properties}
\begin{itemize}[nosep]
\item \textbf{Descent.}  Inequality (3) implies $\Phi(x_{k+1})\le\Phi(x_{k})$ for all $k$ (see proof).
\item \textbf{Hyperparameter Sensitivity.} The bound scales as $1/(\eta T)$; a larger stepsize (up to $1/L$) yields a tighter constant but must respect the $L$‑smoothness condition.
\item \textbf{Comparison.} Compared with plain gradient descent on $f$ alone, PGM inherits the same $O(1/\sqrt{T})$ stationarity rate while handling nonsmooth $g$ via the proximal step.
\end{itemize}

\section*{Discussion}
The result does **not** require convexity of $f$; only $L$‑smoothness is needed. The proximal operator’s \emph{non‑expansiveness} (see Lemma~\ref{lem:nonexp}) is the crucial tool that allows us to bound the distance between successive iterates and to obtain the descent inequality (5) used in the theorem.

\newpage

\section*{Preliminaries (Definitions and Lemmas)}
\begin{lemma}[Non‑expansiveness of the proximal operator]\label{lem:nonexp}
For any $u,v\in\R^{d}$ and any $\eta>0$,
\[
\bigl\|\prox_{\eta g}(u)-\prox_{\eta g}(v)\bigr\|
\;\le\;
\|u-v\|.
\]
\end{lemma}
\begin{proof}
The proximal operator is the resolvent of the subdifferential $\partial g$, which is firmly non‑expansive; the inequality follows directly from the definition of firm non‑expansiveness. \hfill $\square$
\end{proof}

\begin{lemma}[Descent Lemma for $L$‑smooth functions]\label{lem:descent}
If $f$ is $L$‑smooth, then for any $x$ and any $d\in\R^{d}$,
\[
f(x+d)\le f(x)+\langle\grad f(x),d\rangle+\frac{L}{2}\|d\|^{2}. \tag{5}
\]
\end{lemma}
\begin{proof}
Standard consequence of $L$‑smoothness; see \eqref{2}. \hfill $\square$
\end{proof}

\section*{Main Proof (Step‑by‑Step Derivation)}
We prove Theorem~\ref{thm:main}.  Throughout the proof we keep the notation  

\[
v_{k}=x_{k}-\eta\grad f(x_{k}),\qquad
x_{k+1}= \prox_{\eta g}(v_{k}),
\qquad
G_{\eta}(x_{k})=\frac{1}{\eta}(x_{k}-x_{k+1}).
\]

\begin{enumerate}[label={[\textbf{Step \arabic*}]}, leftmargin=*, nosep]
\item \textbf{Optimality condition of the proximal step.}  
By definition of $\prox_{\eta g}$, $x_{k+1}$ satisfies  

\[
0\in \partial g(x_{k+1})+\frac{1}{\eta}(x_{k+1}-v_{k})
   =\partial g(x_{k+1})+\grad f(x_{k})+G_{\eta}(x_{k}). \tag{6}
\]

Hence  

\[
-\grad f(x_{k})-G_{\eta}(x_{k})\in\partial g(x_{k+1}). \tag{7}
\]

\item \textbf{Apply the descent lemma to $f$.}  
Using Lemma~\ref{lem:descent} with $d:=x_{k+1}-x_{k}$,

\[
f(x_{k+1})\le f(x_{k})+\langle\grad f(x_{k}),x_{k+1}-x_{k}\rangle
               +\frac{L}{2}\|x_{k+1}-x_{k}\|^{2}. \tag{8}
\]

Since $x_{k+1}-x_{k}= -\eta G_{\eta}(x_{k})$, we rewrite (8) as  

\[
f(x_{k+1})\le f(x_{k})-\eta\langle\grad f(x_{k}),G_{\eta}(x_{k})\rangle
               +\frac{L\eta^{2}}{2}\|G_{\eta}(x_{k})\|^{2}. \tag{9}
\]

\item \textbf{Use convexity of $g$ via subgradient inequality.}  
From (7) and the subgradient inequality for $g$ we have for any $u$,

\[
g(u)\ge g(x_{k+1})+\langle -\grad f(x_{k})-G_{\eta}(x_{k}),\,u-x_{k+1}\rangle.
\]

Setting $u:=x_{k}$ yields  

\[
g(x_{k})\ge g(x_{k+1})+\langle -\grad f(x_{k})-G_{\eta}(x_{k}),\,x_{k}-x_{k+1}\rangle.
\]

Because $x_{k}-x_{k+1}= \eta G_{\eta}(x_{k})$, we obtain  

\[
g(x_{k})\ge g(x_{k+1})-\eta\langle\grad f(x_{k}),G_{\eta}(x_{k})\rangle
               -\eta\|G_{\eta}(x_{k})\|^{2}. \tag{10}
\]

\item \textbf{Combine (9) and (10).}  
Adding the two inequalities yields a bound on the composite objective:

\[
\begin{aligned}
\Phi(x_{k+1})&=f(x_{k+1})+g(x_{k+1})\\
&\le f(x_{k})+g(x_{k})
   -\eta\langle\grad f(x_{k}),G_{\eta}(x_{k})\rangle
   +\frac{L\eta^{2}}{2}\|G_{\eta}(x_{k})\|^{2}\\
&\qquad\;\; -\eta\langle\grad f(x_{k}),G_{\eta}(x_{k})\rangle
   -\eta\|G_{\eta}(x_{k})\|^{2}\\[2mm]
&= \Phi(x_{k})
   -2\eta\langle\grad f(x_{k}),G_{\eta}(x_{k})\rangle
   -\eta\|G_{\eta}(x_{k})\|^{2}
   +\frac{L\eta^{2}}{2}\|G_{\eta}(x_{k})\|^{2}. \tag{11}
\end{aligned}
\]

\item \textbf{Bound the inner product term.}  
By Cauchy–Schwarz and the elementary inequality $2ab\le a^{2}+b^{2}$,

\[
2\eta\langle\grad f(x_{k}),G_{\eta}(x_{k})\rangle
\le 2\eta\|\grad f(x_{k})\|\|G_{\eta}(x_{k})\|
\le \eta\|\grad f(x_{k})\|^{2}
   +\eta\|G_{\eta}(x_{k})\|^{2}. \tag{12}
\]

Insert (12) into (11) to obtain  

\[
\Phi(x_{k+1})\le \Phi(x_{k})
   -\eta\|\grad f(x_{k})\|^{2}
   -\Bigl(\eta-\frac{L\eta^{2}}{2}\Bigr)\|G_{\eta}(x_{k})\|^{2}. \tag{13}
\]

\item \textbf{Use the stepsize condition.}  
Assumption~\ref{ass:step} gives $\eta\le 1/L$, i.e. $L\eta\le1$. Hence  

\[
\eta-\frac{L\eta^{2}}{2}
   =\eta\Bigl(1-\frac{L\eta}{2}\Bigr)
   \ge \frac{\eta}{2}. \tag{14}
\]

Applying (14) to (13) yields the fundamental descent inequality  

\[
\Phi(x_{k+1})\le \Phi(x_{k})-\frac{\eta}{2}\|G_{\eta}(x_{k})\|^{2}. \tag{15}
\]

\item \textbf{Summation over iterations.}  
Telescoping (15) from $k=0$ to $T-1$ gives  

\[
\Phi(x_{T})\le \Phi(x_{0})-\frac{\eta}{2}\sum_{k=0}^{T-1}\|G_{\eta}(x_{k})\|^{2}. \tag{16}
\]

Rearranging and using Assumption~\ref{ass:bound} ($\Phi(x_{T})\ge\Phi^{\inf}$) leads to  

\[
\sum_{k=0}^{T-1}\|G_{\eta}(x_{k})\|^{2}
\le \frac{2}{\eta}\bigl(\Phi(x_{0})-\Phi^{\inf}\bigr). \tag{17}
\]

Dividing by $T$ yields exactly inequality \eqref{3}. \hfill $\square$
\end{enumerate}

\section*{Rate Derivation (With Constants)}
From \eqref{3} we have  

\[
\min_{0\le k\le T-1}\|G_{\eta}(x_{k})\|
\le\frac{1}{T}\sum_{k=0}^{T-1}\|G_{\eta}(x_{k})\|
\le\sqrt{\frac{2}{\eta T}\bigl(\Phi(x_{0})-\Phi^{\inf}\bigr)}.
\]

Thus the algorithm attains an $\epsilon$‑stationary point (i.e. a point with $\|G_{\eta}(x)\|\le\epsilon$) after at most  

\[
T\ge \frac{2}{\eta\epsilon^{2}}\bigl(\Phi(x_{0})-\Phi^{\inf}\bigr)
\]

iterations.  The rate $O(1/\sqrt{T})$ for the norm of the gradient mapping matches the optimal rate for first‑order methods in the non‑convex setting.

\section*{Special Cases}
\begin{itemize}[nosep]
\item \textbf{Convex $f$.}  If $f$ is convex, the same bound (15) holds, and one can further show that $\Phi(x_{k})$ converges to the optimal value at rate $O(1/k)$ when $\eta=1/L$.
\item \textbf{Pure gradient descent ($g\equiv0$).}  Then $\prox_{\eta g}$ is the identity, $G_{\eta}(x)=\grad f(x)$, and (15) reduces to the classical descent lemma for smooth non‑convex functions.
\item \textbf{Exact proximal step ($\eta\to0$).}  As $\eta\downarrow0$, $G_{\eta}(x)\approx \grad f(x)+\partial g(x)$, and the bound recovers the standard stationarity condition for the original problem.
\end{itemize}

\end{document}